\documentclass[11pt]{article}

\usepackage[margin=1in]{geometry}
\usepackage{hyperref}
\usepackage{amsmath}
\usepackage{enumitem}
\usepackage{listings}
\usepackage{xcolor}

\hypersetup{
    colorlinks=true,
    linkcolor=blue,
    urlcolor=blue,
    citecolor=blue
}

\title{CS 300 -- Week 3 Resources\\\large Linked Lists: Deletion in a Doubly Linked List}
\author{Jose de Lima}
\date{\today}

\begin{document}
\maketitle

\section*{Overview}
Week 3 shifts from array-backed structures to \textbf{linked lists}, where
each element (a \emph{node}) owns its own memory and references its neighbors
through pointers. The resource below focuses on one of the trickiest
operations on a \textbf{doubly linked list}: deleting a node, whether it
sits at the head, in the middle, or at the tail.

\section{Doubly Linked List Refresher}

A \textbf{doubly linked list (DLL)} node stores three things:
\begin{itemize}[leftmargin=*]
    \item \texttt{data} -- the payload (an \texttt{int}, a \texttt{Bid}, etc.).
    \item \texttt{prev} -- a pointer to the previous node (or \texttt{nullptr}
          if this is the head).
    \item \texttt{next} -- a pointer to the next node (or \texttt{nullptr} if
          this is the tail).
\end{itemize}

The list itself keeps a \texttt{head} pointer (and usually a \texttt{tail}
pointer for $O(1)$ tail access). Compared to a singly linked list, the extra
\texttt{prev} pointer costs memory but makes backward traversal and node
deletion far easier, because a node can find its predecessor in $O(1)$
without re-walking the list.

\begin{center}
\texttt{NULL <- [prev|data|next] <-> [prev|data|next] <-> [prev|data|next] -> NULL}
\end{center}

\section{Deleting a Node from a Doubly Linked List}
\textbf{Resource:} \href{https://www.youtube.com/watch?v=LvUgew66zOQ}{Delete a node from a Doubly Linked List (YouTube)}

Deletion breaks into three cases depending on where the target node sits in
the list. In each case the goal is the same: detach the node from its
neighbors, update the surrounding pointers so the list remains well-formed,
and free the node's memory.

\subsection{Case 1 -- Delete the Head Node}
\begin{enumerate}[leftmargin=*]
    \item Let \texttt{temp = head}.
    \item Advance the head: \texttt{head = head->next}.
    \item If the new \texttt{head} is not \texttt{nullptr}, set
          \texttt{head->prev = nullptr} (the new head has no predecessor).
    \item If the new \texttt{head} is \texttt{nullptr}, the list is now empty,
          so also clear \texttt{tail}.
    \item \texttt{delete temp;}
\end{enumerate}

\subsection{Case 2 -- Delete a Middle Node}
Assume \texttt{cur} points at the node to remove and it is not the head or
tail.
\begin{enumerate}[leftmargin=*]
    \item \texttt{cur->prev->next = cur->next;}
    \item \texttt{cur->next->prev = cur->prev;}
    \item \texttt{delete cur;}
\end{enumerate}
This is the case where DLLs shine -- no second traversal is needed to find
the predecessor, because \texttt{cur->prev} already points at it.

\subsection{Case 3 -- Delete the Tail Node}
\begin{enumerate}[leftmargin=*]
    \item Let \texttt{temp = tail}.
    \item Move tail back: \texttt{tail = tail->prev}.
    \item If the new \texttt{tail} is not \texttt{nullptr}, set
          \texttt{tail->next = nullptr}.
    \item If the new \texttt{tail} is \texttt{nullptr}, also clear
          \texttt{head} (list is empty).
    \item \texttt{delete temp;}
\end{enumerate}

\subsection{Pseudocode (all three cases combined)}
\begin{lstlisting}[basicstyle=\ttfamily\small,frame=single]
function deleteNode(list, cur):
    if cur == nullptr: return

    if cur->prev != nullptr:
        cur->prev->next = cur->next
    else:
        list.head = cur->next      # cur was the head

    if cur->next != nullptr:
        cur->next->prev = cur->prev
    else:
        list.tail = cur->prev      # cur was the tail

    delete cur
\end{lstlisting}

\subsection{Edge Cases to Watch}
\begin{itemize}[leftmargin=*]
    \item \textbf{Empty list:} nothing to delete -- return early.
    \item \textbf{Single-node list:} the node is \emph{both} head and tail;
          after deletion, both must be set to \texttt{nullptr}.
    \item \textbf{Dangling pointers:} always \texttt{delete} the removed node
          and stop using references to it afterward.
    \item \textbf{Double-free:} make sure no other structure still owns the
          node before freeing.
\end{itemize}

\subsection{Complexity}
\begin{itemize}[leftmargin=*]
    \item \textbf{If you already hold a pointer to the node:} $O(1)$ time,
          $O(1)$ space -- this is the doubly linked list's signature win
          over a singly linked list.
    \item \textbf{If you must first find the node by value/index:} $O(n)$
          time for the search, plus $O(1)$ for the unlink.
\end{itemize}

\section{Singly vs. Doubly Linked List -- Deletion}
\begin{center}
\begin{tabular}{l|l|l}
\textbf{Property}               & \textbf{Singly} & \textbf{Doubly} \\ \hline
Pointers per node               & 1 (\texttt{next})                & 2 (\texttt{prev}, \texttt{next}) \\
Delete with node pointer only   & $O(n)$ (need predecessor)        & $O(1)$                            \\
Backward traversal              & Not supported in $O(1)$          & $O(1)$ per step                   \\
Memory overhead                 & Lower                             & Higher (extra pointer)            \\
Implementation complexity       & Simpler                           & More pointer bookkeeping          \\
\end{tabular}
\end{center}

\section*{Key Takeaways}
\begin{itemize}[leftmargin=*]
    \item Doubly linked lists trade one extra pointer per node for $O(1)$
          deletion when you already hold the node.
    \item Every DLL deletion is fundamentally the same three-step dance:
          relink the predecessor, relink the successor, free the node --
          with special handling when either neighbor is \texttt{nullptr}.
    \item Always reason about the empty-list and single-node cases
          explicitly; most linked-list bugs live there.
\end{itemize}

\section*{Links Summary}
\begin{itemize}[leftmargin=*]
    \item Delete a Node from a Doubly Linked List: \url{https://www.youtube.com/watch?v=LvUgew66zOQ}
\end{itemize}

\end{document}
